\bookchapter{Kapitel 63: Views, State und Navigation}

\chapterintro{Einseitige Spielereien kannst du bereits. Dieses Kapitel löst die zwei Fragen des „Erwachsenwerdens“: wie du zwischen mehreren Seiten hin- und herspringst (Navigation) und wie Daten zwischen den Seiten fließen (State-Management). Nach diesem Kapitel baust du das klassische Mehrseiten-Gerüst „Liste Detail Bearbeiten“.}

\booksection{Das Datenfluss-Panorama des State-Managements}

SwiftUI hat mehr als nur @State als State-Werkzeug. Sieh dir zuerst diese Übersicht an, um zu verstehen, „wann du welches einsetzt“:

\begin{center}
\begin{tabularx}{\linewidth}{|X|X|X|X|}
\hline
\textbf{Werkzeug} & \textbf{Zuständig für} & \textbf{Wer es nutzt} & \textbf{Typischer Einsatzfall} \tabularnewline
\hline
@State & Seiteneigener, temporärer Zustand & Das aktuelle View & Zähler, Textfeldeingaben \tabularnewline
@Binding & Zustand von jemand anderem ausleihen & Kind-View ändert Eltern-View & Eigene Komponenten \tabularnewline
@StateObject / @Observable & Seitenübergreifend geteilte Logik & ViewModel & Datenladen, Geschäftslogik \tabularnewline
@EnvironmentObject / .environment & Global geteilter Zustand & Die ganze App & Anmeldestatus, Theme-Einstellungen \tabularnewline
\hline
\end{tabularx}
\end{center}

\begin{center}\small\color{gray}Überblick über SwiftUIs State-Werkzeuge (2026: @Observable ist längst Standard)\end{center}

\begin{lstlisting}[language=Swift]
// der moderne Weg ab iOS 17+: das @Observable-Makro (für neue Projekte empfohlen)
import Observation

@Observable
class CounterModel {                 // das View Model
    var count = 0
    func increment() { count += 1 }
}

struct CounterView: View {
    @State private var model = CounterModel()   // hält das Modell

    var body: some View {
        VStack {
            Text("Du hast \(model.count) Mal geklickt")
            Button("Klick mich") { model.increment() }
        }
    }
}
\end{lstlisting}

Knallharter Rat für Einsteiger: Seiteneigen @State; seitenübergreifend oder mit Logik @Observable-Klasse. @StateObject/@ObservableObject aus älteren Tutorials ist die Schreibweise der Vorgängergeneration – lesen können genügt, für neuen Code gilt das oben Gesagte.

\booksection{NavigationStack: der „Stapel“ für Seitenwechsel}

Ab iOS 16 übernimmt NavigationStack die Mehrseiten-Navigation – wie ein Stapel Teller: einer kommt oben drauf, und beim Zurückgehen wird jeder einzeln wieder abgenommen:

\begin{lstlisting}[language=Swift]
struct HomeView: View {
    var body: some View {
        NavigationStack {
            List {
                NavigationLink("Max' Ausgabenbuch") {
                    DetailView(name: "Max' Ausgabenbuch")   // Zielseite
                }
                NavigationLink("Lenas Ausgabenbuch") {
                    DetailView(name: "Lenas Ausgabenbuch")
                }
            }
            .navigationTitle("Liste der Ausgabenbücher")   // große Überschrift oben
        }
    }
}

struct DetailView: View {
    let name: String
    var body: some View {
        Text("Details zu \(name)")
            .navigationTitle(name)
            .navigationBarTitleDisplayMode(.inline)
    }
}
\end{lstlisting}

NavigationLink ist schlicht ein „antippbarer Sprunglink“; navigationTitle erzeugt automatisch Navigationstitel und Zurück-Button. Für programmatisches Navigieren (etwa nach „Speichern“ automatisch zurück) schließt du die aktuelle Seite mit @Environment(\textbackslash{}.dismiss) – mit dieser Kombination hast du 90 \% aller Navigationsbedürfnisse abgedeckt.

\booksection{Dialoge und untere Reiter: sheet und TabView}

Neben dem „Seitenschieben“ gibt es zwei klassische Interaktionsformen:

\begin{lstlisting}[language=Swift]
// 1. Sheet von unten, halber oder ganzer Bildschirm (der Liebling von Eingabeformularen)
struct HomeView: View {
    @State private var showAdd = false
    var body: some View {
        Button("Ausgabe erfassen") { showAdd = true }
            .sheet(isPresented: $showAdd) {
                AddRecordView()     // die sich öffnende Erfassungsansicht
            }
    }
}

// 2. untere Tab-Leiste (das "Skelett" der App)
struct MainTabView: View {
    var body: some View {
        TabView {
            RecordListView()
                .tabItem { Label("Buch", systemImage: "book") }
            StatsView()
                .tabItem { Label("Statistik", systemImage: "chart.pie") }
            SettingsView()
                .tabItem { Label("Einstellungen", systemImage: "gear") }
        }
    }
}
\end{lstlisting}

\begin{tcbinfo}{Tipp: Wann nimmt man was?}
Faustregel: Für den Hauptablauf Seiten mit NavigationStack schieben (Liste Detail), für kurzzeitige Aktionen den sheet-Dialog (Anlegen, Bearbeiten, Bestätigen) und für gleichrangige Hauptbereiche TabView (untere Reiter). Die Kombination aus allen dreien ist das Gerüst der allermeisten Apps.

\end{tcbinfo}

\booksection{List: iOS' Alleskönner-Liste}

Der meistgenutzte und am typischsten „apple-artige“ Baustein unter iOS ist List – mit Trennlinien, Wisch-löschen und Tipp-Highlight ab Werk:

\begin{lstlisting}[language=Swift]
struct RecordListView: View {
    @State private var records = ["Mittagessen 25.5", "U-Bahn 6.0", "Gehalt +8000"]

    var body: some View {
        List {
            ForEach(records, id: \.self) { r in
                Text(r)
            }
            .onDelete { indexSet in          // nach links wischen zum Löschen
                records.remove(atOffsets: indexSet)
            }
        }
    }
}
\end{lstlisting}

Liste + NavigationLink + sheet + TabView – kombiniere diese vier Bausteine, und du kannst bereits das komplette App-Gerüst „Buchungsliste Detail öffnen auf + tippen, Eingabe-Dialog erscheint“ bauen. Genau so entsteht das Haushaltsbuch-Projekt im nächsten Kapitel.

\begin{tcbpractice}{Jetzt bist du dran: Baue ein „Drei-Seiten-Gerüst“}
\booksection{Der App-Lebenszyklus: Vordergrund und Hintergrund}

Ähnlich wie der Activity-Lebenszyklus bei Android durchläuft auch eine iOS-App Zustandswechsel von „Geburt, Vordergrund, Hintergrund, Tod“. In SwiftUI beobachtest du das über den Environment-Wert scenePhase:

\begin{lstlisting}[language=Swift]
@main
struct MyApp: App {
    @Environment(\.scenePhase) private var phase

    var body: some Scene {
        WindowGroup {
            ContentView()
        }
        .onChange(of: phase) { _, newPhase in
            switch newPhase {
            case .active:    print("wieder vorn: Aktualisierung fortsetzen")
            case .inactive:  print("im Übergang (Anruf oder Moment des Wechsels nach hinten)")
            case .background: print("in den Hintergrund: Entwurf sichern")
            @unknown default: break
            }
        }
    }
}
\end{lstlisting}

Typische Einsätze: Daten auffrischen, wenn die App wieder in den Vordergrund kommt; Entwürfe speichern, wenn sie in den Hintergrund geht. iOS geht mit Hintergrundbetrieb streng ins Gericht – normale Apps werden bald nach dem Wechsel in den Hintergrund angehalten (eingefroren). Rechne nicht damit, heimlich Schwerstarbeit im Hintergrund zu leisten; für zeitgesteuerte Aufgaben nutze die lokalen Benachrichtigungen aus Kapitel 68, keinen „Dauerbetrieb im Hintergrund“.

Baue eine Fake-Daten-Version des Haushaltsbuch-Gerüsts: TabView mit drei Reitern (Buchungen/Statistik/Einstellungen) die Buchungsseite zeigt Fake-Daten in einer List, ein Tipp auf eine Zeile springt zum Detail (NavigationStack) der “+”-Button oben rechts öffnet den Eingabe-Dialog als sheet. Kein Speichern, kein Netz – reine Gerüstübung. Steht dieses Gerüst, ist es in Kapitel 65 nur noch ein Katzensprung, echte Daten einzufüllen.

\end{tcbpractice}

\begin{tcbsummary}{Kapitelfazit}
\begin{itemize}
  \item Die vier State-Werkzeuge: @State (seiteneigen), @Binding (ausgeliehen), @Observable (Logik), .environment (global)
  \item In neuen Projekten schreibst du ViewModels mit dem @Observable-Makro; @ObservableObject aus alten Tutorials: lesen können genügt
  \item NavigationStack schiebt Seiten; @Environment(\textbackslash{}.dismiss) ermöglicht programmatisches Zurück
  \item Hauptablauf: Seiten schieben; kurzzeitige Aktionen: sheet; gleichrangige Hauptbereiche: TabView
  \item List ist die Alltagsliste: Trennlinien und Wisch-löschen ab Werk
\end{itemize}

\end{tcbsummary}
